\section{Additional Related Works}
\paragraph{Data Reuse in Synthetic Setting.} Besides the real-world LLM pre-training regime, many works also reported the improvement of data reusing under synthetic settings empirically~\citep{charton2024emergent,kazdan2024collapse} or theoretically~\citep{zucchet2025emergence,dandi2024benefitsreusingbatchesgradient,arnaboldi2025repetitaiuvantdatarepetition}. 

\paragraph{Empirical Findings on Scaling Laws.} Scaling laws reveal the relationships between large-scale model training loss and various factors such as model size, data size, and compute budget. These laws were initially observed by \citet{hestness2017deep}, but gained significant influence through the work of \citet{kaplan2020scaling}, and have since been further developed in a series of studies \citep{henighan2020scaling, hoffmann2022training, zhai2022scaling, kadra2023power, aghajanyan2023scaling, muennighoff2023scaling, bi2024deepseek, shuai2024scaling, kumar2024scaling, tissue2024scaling, luo2025multi}. %Notably, \citet{hoffmann2022training} introduced the Chinchilla Law, which proposes an optimal scaling strategy under a fixed compute budget. Additionally, \citet{tissue2024scaling} and \citet{luo2025multi} explored scaling laws with learning rate schedules (LRSs) based on empirical observations. 
Notably, \citet{muennighoff2023scaling} further refined these models by incorporating the number of training epochs into a more complex scaling law, which empirically describes the effect of data reuse. In our work, we provide a theoretical analysis of how the effective reuse rate $E(K,N)$ relies on the epoch number $K$ and fresh data size $N$, highlighting the role of $N$ in the scaling behavior of $E(K,N)$, a factor that was overlooked in \citet{muennighoff2023scaling}. 

% \paragraph{Theory of Scaling laws.}
% There are many existing works that interpreted the mechanism behind the scaling laws ... . \citet{sharma2020neuralscalinglawdimension} considered a toy model and showed that $\cL \propto \frac{1}{N^{\frac{4}{d}}}$, where $d$ is the intrinsic dimension. \citet{Bahri_2024} and \citet{bordelon2024dynamicalmodelneuralscaling} both analyzed a linear teacher-student model setting with a power-decay parameter prior. They proposed theories for scaling laws based on statistical heruistics. \citet{bordelon2025featurelearningimproveneural} developed a solvable model of neural scaling laws and defined regimes according to task difficulties. Meanwhile, other works considered more complicated settings and gave more rigorous mathematical proofs. \citet{jain2024scalinglawslearningreal} studied how data from other sources affect the test error of models and their scaling laws. \citet{paquette202543phasescomputeoptimalneural} identified different phases for a PLRF model in the compute-limited, infinite-data scaling law regime, and they discovered different scaling exponents in different phases. \citet{lin2024scaling} proved a scaling law using a PLRF model with an online SGD scheme with a geometrically decaying stepsize non-asymptotically. However, these scaling law papers failed to consider the impact of data reuse. In contrast, we proved matching upper and lower risk bounds in two settings for multi-epoch training.
% \paragraph{Theory of Multi-Pass SGD.}
% Shuffling is a common practice in machine learning, however, theoretical literature on SGD with shuffling is limited. 

%\haodong{emphasize that the linear regression setting is attracting more and more interest}
\paragraph{Theoretical Explanations for Scaling Laws.} A series of studies \citep{sharma2020neuralscalinglawdimension,hutter2021learning,maloney2022solvable,wei2022more,jain2024scalinglawslearningreal,michaud2024quantization,nam2024exactly,atanasov2024scaling,dohmatob2024tale,bahri2024explaining,bordelon2024dynamical,lin2024scaling,paquette202543phasescomputeoptimalneural,bordelon2024feature,zhang2024does,ferbach2025dimension,li2025functional,li2026muon,li2026optimal,wang2026fast} have sought to theoretically explain scaling laws from various perspectives. Among these, recent works \citep{bordelon2024dynamical,paquette202543phasescomputeoptimalneural,lin2024scaling,bordelon2024feature} have analyzed scaling laws by tracking the training dynamics of SGD through linear regression setup. Specifically, \citet{bordelon2024dynamical} investigated a full-batch gradient flow setup, while \citet{paquette202543phasescomputeoptimalneural} and \citet{bordelon2024feature} focused on online SGD with a sufficiently small constant learning rate. Additionally, \citet{lin2024scaling} studied a geometric decaying learning rate schedule (LRS) \citep{ge2019step,wu2022iterateriskboundssgd}. Recently, \citet{li2025functional} proposed a functional scaling law that characterizes the loss dynamics for general LRSs. However, these scaling law studies did not account for the impact of data reuse. In contrast, our work examines the scaling behavior of multi-epoch SGD training within the context of a linear regression setup. 

\paragraph{SGD Analysis in Linear Regression.} %\haodong{mention some scaling law papers} 
The analysis of SGD in linear regression has been extensively studied over the years, encompassing both one-pass and multi-epoch SGD. In the context of one-pass SGD, \citet{zou2021benign,meterez2025simplified} considered an SGD procedure with a constant step size and averaged iterates, offering a sharp risk bound in terms of the eigenvalues of the covariance matrix. \citet{pmlr-v139-gurbuzbalaban21a} examined one-pass SGD with batch size and proved that the distribution of
the SGD iterates will converge to a heavy-tailed
stationary distribution. \citet{zou2022benefitsimplicitregularizationsgd} compared the performance of SGD in the absence of ridge regression. \citet{wu2022iterateriskboundssgd} and \citet{wu2022powerlimitationpretrainingfinetuninglinear} studied SGD in linear regression under covariate shift. \citet{xia2024data} considered SGD updates with noisy gradient and analyzed the perfect deleted point problem. \citet{li2025understanding} considered SGD with exponential moving average in the linear regression setting. For multi-epoch SGD, \citet{lin2019optimalratesmultipassstochastic} examined a scenario in which gradients are sampled uniformly at random and mini-batches are allowed. They analyzed the effects of mini-batch size, number of epochs, and learning rate, carefully combining these parameters to achieve the optimal convergence rate. \citet{pillaudvivien2018statisticaloptimalitystochasticgradient} showed that while single-pass averaged SGD is optimal for a certain class of "easy" problems, multiple passes are required to achieve optimal prediction performance on a different class of "hard" problems, provided that an appropriate step size is chosen. In contrast to the matching upper and lower bounds derived by our theory, however, all the above works were only able to derive an upper bound for the loss.